%% kapa: exam=1 round=37 year=2026 period=2 form=A subject=accounting subject_ko=회계학 seq=26 number=66
%% source: Q-Net 감정평가사 기출문제 (cst00309) artlSeq=5258924 file=raw/1008/5258924/2026년 제37회 감정평가사 1차 2교시.pdf
%% conversion: kordoc 4.24.51 -> tools/gen1.py (auto); review=proofread-sonnet; flags=none
\begin{question}{66}
(주)감평은 20×1년 초 임대목적 투자부동산인 상가건물을 ￦1,000,000에 취득(내용연수 20년, 잔존가치 ￦0, 정액법 상각)하고 원가모형을 적용하여 회계처리하였다. 20×2년 초부터 동 자산에 대해 원가모형 대신 공정가치모형을 적용하기로 하였으며, 이러한 회계변경은 정당한 변경에 해당한다. 동 건물의 공정가치는 다음과 같으며, (주)감평의 20×1년 말 이익잉여금은 ￦2,000,000이었다. 투자부동산에 대한 회계처리 반영 전 20×2년 당기순이익이 ￦500,000일 경우, 상기 회계변경 반영 후 (주)감평의 20×2년 말 이익잉여금은? (단, 상기 건물 이외의 투자부동산은 보유하지 않으며, 이익잉여금 처분은 없다.)
\par
{\small\centering
\begin{tabular}{|c|c|c|}\hline
연도 & 20×1년 말 & 20×2년 말 \\\hline
공정가치 & ￦800,000 & ￦700,000 \\\hline
\end{tabular}\par}
\begin{choices}
  \item ￦1,850,000
  \item ￦2,000,000
  \item ￦2,200,000
  \item ￦2,250,000
  \item ￦2,500,000
\end{choices}
\end{question}
\answer{4}
